\subsection{Summary}
\label{sec:summary}

The exploration of the literature in the related thematic areas, highlights the level of research undertaken to realise the smart grid and the challenges faced in doing so. In particular,
resiliency of the grid is referenced frequently as a key goal, with FLISR providing an effective method of achieving this objective. In regards to FLISR, there are many applications in the research
that utilise various methods of fault identification and resolution in terms of the algorithms applied, communication protocols, standards and methods of control, such as fully autonomous
FLISR as well as those that require human intervention. A key concern identified in the literature is the need to ensure such grid resiliency solutions are in themselves fault-tolerant and therefore not dependent on centralised
systems or core telecommunication networks. In reviewing the research, distributed FLISR solutions appear to be a means of resolving this limitation.
\\
In realising a distributed FLISR solution edge computing emerges as an ideal paradigm to implement this functionality.
By removing the need for large data processing and decision making functionality from cloud-based systems and moving subsets of this logic to the edge of the network, grid resiliency solutions,
such as distributed FLISR, should benefit from reduced latency and dependency on core systems in order to operate effectively.
Indeed, the literature describes how edge computing may be a key component in realising the smart grid, due to the amount of data generated and computation required to effectively manage the
energy distribution system. In combination with this, is the difficultly in effecting large scale changes on the core distribution grid due to expense and dealing with legacy components. Therefore, enabling
more computation, data collection and control at the edge of the network may afford smart grid functionality, implementation of concepts such as the VPP, microgrids and prosumer activities.

\clearpage
% subsection summary (end)
